%!TEX root = ../../thesis.tex

\section{Actions}\label{sec:mdp-actions}

EfficientZero requires a discrete action space of tractable size, usually on the
order of tens to hundreds of actions.
In contrast to \gls{rl} algorithms like Q-Learning that can deal well with probabilistic
transitions, EfficientZero requires actions to have a deterministic effect to allow proper planning.
This rules out the direct use of probabilistic mutation operators like random
bit flips, bit insertions or structural mutations like random splicing.
Designing a deterministic set of mutation operators is not easy since the effect
of an action must be compatible with the observation model.

Consider for example a seed $s$ that has been assigned a power of
$4\leftarrow\textsc{AssignEnergy}(s)$ and on which a
mutation operator $m$ generates four new seeds $m(s,4)=\{s_0,s_1,s_2,s_3\}$
that must be reconciled into a single observation since the mutation operator is modeled
as a single action, even if it generates multiple test cases.
Reconciliation strategies and additional mutation operators are outlined in~\cref{chap:future-work}.

\subsection{Input-level Actions}\label{sec:input-level-actions}

One way to model actions is as actual \gls{sut} inputs.\index{action space!input-level}
This is a viable approach for \glspl{sut} that consume their inputs byte-by-byte.
The learned policy $\pi$ in this case represents a generative model for inputs
that conform to the \gls{sut}'s grammar.

\subsection{Corpus-level Actions}\label{sec:corpus-level-actions}

An action model that is comparable to traditional \gls{cgf} approaches is that of
corpus-level actions.\index{action space!corpus-level}
In general, we keep track of a seed corpus $\mathcal{S}$ that is evolved during an episode.
A single action then decides three things: whether to retain the previous seed,
which seed to select next, and how to mutate the current seed.
In the simplified corpus setting, we keep track of only a single seed $s$ of fixed length $l$.
Actions then represent a binary retain/discard decision, as well as mutations on $s$.
Further, we include $s$ in the observations to allow the agent to base its decisions
on the coverage as well as the input that produced this coverage.
In theory, coverage is fully determined by $s$, which is why~\textcite{bottinger_deep_2018}
chose to set the state space equal to the input space.
We claim that keeping coverage information alongside the seed is advantageous.

Classically, mutations are implemented by selecting a location and operation.
We choose to keep the location selection, but replace the operator selection with
the direct selection of a byte value to write to the selected location.
This approach has two advantages.
First, we rid ourselves of any probabilistic mutation operators.
While these are usually an integral part of mutational fuzzers, their probabilistic
nature is ill-suited to MuZero-style architectures since it makes planning more unpredictable.
Second, we do not need a handcrafted choice of mutation operators.
We expect that the agent can learn a suitable set of operators itself,
especially since the current seed state $s$ is included in the
observations.
